\section{来源审计与课堂边界：这是一堂 2023 年的导论课}

本讲最需要先处理的不是公式，而是证据边界。视频前 10 分钟由 CS25 课程团队介绍系列目标，10:14 之后才进入 Andrej Karpathy 的主体讲座；课堂日期是 2023 年 1 月 10 日，Stanford Online 的上传日期则是 2023 年 5 月 19 日。旧稿把上传日期写成 2023 年 2 月 23 日，并加入大量课堂没有讲到的 RLHF 治理、生产部署、drift detection 与团队流程，因此本次不是增量润色，而是按官方视频、人工字幕和一手论文完整重写。

\lecturefigure{slide-01-course-title.jpg}{CS25 Transformers United V2：课程系列的正式开场。}{Stanford Online 官方视频 00:00:05--00:00:52。}

这张标题页只承担课程身份识别，却也给出一个重要阅读提示：这里的 ``Transformers'' 不是单一 NLP 模型，而是一个跨领域 speaker series 的共同语言。讲义后面会把这种跨领域能力拆成三层：输入如何被表示成元素序列、元素之间如何通过 attention 通信，以及残差网络如何把通信结果转化为可训练的计算。

\lecturefigure{slide-02-course-content-logistics.jpg}{课程内容与组织说明：以跨领域嘉宾讲座理解 Transformer。}{Stanford Online 官方视频 00:00:52--00:02:53。}

课程团队并没有承诺一堂课覆盖全部细节，而是用本讲建立共同词汇，再让后续嘉宾展示 Transformer 在不同领域如何被改造。读这页时应把 ``共同架构'' 与 ``相同系统'' 分开：不同课程讲座会共享 attention、residual、token/element representation 等骨架，但数据接口、目标函数、mask 和 inductive bias 仍然高度领域化。

\begin{importantbox}{本讲的证据合同}
课堂中的历史回忆、工程直觉和未来判断会标成“课堂提示”或“讲者判断”；论文中的机制会给出公式；2023 年之后的事实不会倒灌进讲者当时的叙述。没有公开独立 slide PDF，因此正文中的 61 张图均是从官方 1080p 视频人工复核出的教学状态，不冒充官方 deck 页码。
\end{importantbox}

\subsection{课程希望读者带走什么}

本节先把目标压缩成三个问题：Transformer 怎样工作，为什么能走出 NLP，以及哪些性质使它成为一个可持续扩展的研究平台。后文每个技术细节都应服务于至少一个问题；如果一个概念不能说明通信、计算、优化或输入接口，它就不应成为孤立名词。

\lecturefigure{slide-03-learning-goals.jpg}{三项目标：机制、跨领域应用与研究前沿。}{Stanford Online 官方视频 00:02:54--00:03:17。}

\begin{knowledgebox}{读图：三项目标不是三门互不相干的课}
第一项“how Transformers work”提供可执行机制；第二项“beyond NLP”检验机制是否可迁移；第三项“new directions”追问现有机制在哪里失效。正确学习顺序是先理解一个最小 decoder，再观察输入与 connectivity 如何变化，最后讨论 context、memory、cost 与 controllability 等边界，而不是先背应用名单。
\end{knowledgebox}

\lecturefigure{slide-04-transformers-introduction.jpg}{导论主题：从 attention 时间线进入 Transformer。}{Stanford Online 官方视频 00:03:17--00:03:35。}

这页看似只是章节标题，实际把后面的叙事定为“问题链”而非“组件表”。Attention 不是突然出现的模块：它先解决 encoder 把整句压进一个向量的瓶颈，随后被提升为主要通信机制，再与位置、残差、归一化和 MLP 组成完整架构。理解这条链，才能知道每个组件为何存在。

\lecturefigure{slide-05-attention-timeline.jpg}{Attention 时间线：从早期序列模型到 2017 年 Transformer。}{Stanford Online 官方视频 00:03:35--00:05:01。}

\begin{knowledgebox}{读图：时间线应按“瓶颈被怎样消除”阅读}
早期 RNN/LSTM 通过 recurrent state 处理变长序列，却要求信息沿时间逐步传递；seq2seq 又把源句压到固定 context vector；Bahdanau attention 让 decoder 每一步回看全部 encoder states；Transformer 进一步去掉 recurrent path，让任意允许连接的 token 在一层内直接交换信息。时间越往后，信息路径越短、并行度越高，但显式 connectivity 与位置表示也越重要。
\end{knowledgebox}

\lecturefigure{slide-06-prehistoric-rnns.jpg}{“史前时代”：RNN/LSTM 主导序列建模，但受串行路径约束。}{Stanford Online 官方视频 00:05:01--00:06:10。}

\term{RNN}（Recurrent Neural Network，循环神经网络）在每个时间步用上一状态与当前输入更新隐藏状态；\term{LSTM}（Long Short-Term Memory）通过门结构改善长期依赖与梯度传播。它们能表示变长序列，却把训练计算图拉成长而细的链：第 $t$ 个位置依赖前 $t-1$ 步，GPU 很难把时间维完全并行化，监督信号回到早期输入也要经过许多跳。

\lecturefigure{slide-07-where-we-were-2021.jpg}{2021 年的课程视角：长序列、AlphaFold、offline RL、few-shot 与多模态开始汇合。}{Stanford Online 官方视频 00:06:10--00:06:51。}

这页是历史快照，不是 2026 年能力清单。课程团队强调的转折是：Transformer 已经从文本生成扩展到蛋白质、离线强化学习和图文生成，并显露 zero-shot/few-shot generalization。读者应比较的不是今天的绝对性能，而是研究者当时已经看见“一种骨架覆盖多种数据结构”的趋势。

\lecturefigure{slide-08-where-we-are-2022.jpg}{2022 年的扩展：音频、艺术、推理、human feedback 与 diffusion。}{Stanford Online 官方视频 00:06:51--00:07:30。}

课堂把 ChatGPT、human feedback、multimodal generation 与 diffusion models 放在同一时代背景中，但不能由此推出它们都由同一目标训练。Transformer 是序列处理骨架；RLHF 是把人类偏好注入模型行为的训练流程；diffusion 是另一类生成过程。把这些层次混在一起，正是旧稿出现大量无来源治理内容的根本原因。

\lecturefigure{slide-09-future.jpg}{2023 年初的未来问题：视频、长序列、agent、领域模型、外部记忆与更低计算成本。}{Stanford Online 官方视频 00:07:30--00:10:12。}

\begin{warningbox}{读图时必须保留时间戳}
页中的 4,000-token 量级、ChatGPT 无长期记忆、视频模型仍待突破等陈述，是 2023 年 1 月的课堂观察。它们的教学价值在于识别结构性问题：attention 对长度的二次代价、有限 context、随机输出的可控性、外部 memory 与 domain specialization；不能把具体产品状态当成 2026 年事实。
\end{warningbox}

\subsection{Karpathy 的讲座从哪里开始}

课程团队完成背景介绍后，Karpathy 才正式接手。他说自己虽然住得近、可以走来上课，却为 slides 花了十小时，并因时间有限删掉了两个原计划主题。这一口头说明很重要：本讲是刻意压缩的 conceptual tour，后文不应假装它完整覆盖位置编码变体、训练系统、对齐或推理服务。

\lecturefigure{slide-10-karpathy-title.jpg}{Andrej Karpathy 的主体讲座标题页。}{Stanford Online 官方视频 00:10:14--00:10:39。}

\teachervoice{Karpathy 明确说他会简化讲座并跳过部分主题。讲义因此选择“历史问题链 + attention 图解释 + nanoGPT 端到端实现 + 跨模态与 in-context learning”四条主线，而不是向课堂外扩写通用 LLM 全栈。}

\subsection{本章小结}

本章建立了时间、说话人和证据边界：前十分钟属于课程导论，主体讲座始于 10:14；课堂发生在 2023 年 1 月，上传发生在 2023 年 5 月；61 张视觉材料来自官方视频恢复。接下来进入 Karpathy 的核心历史论证：Transformer 的价值首先来自 AI 各领域从手工特征孤岛走向共同可训练架构。

